\documentclass[11pt,a4paper]{report}
\usepackage[utf8]{inputenc}
\usepackage[margin=1in]{geometry}
\usepackage{xcolor}
\usepackage{fancyhdr}
\usepackage{titlesec}
\usepackage{booktabs}
\usepackage{tabularx}
\usepackage{hyperref}
\usepackage{tcolorbox}
\usepackage{enumitem}
\usepackage{listings}
\usepackage{amsmath,amssymb}
\usepackage{graphicx}

% Colors
\definecolor{cdacblue}{RGB}{14, 53, 101}
\definecolor{labcyan}{RGB}{0, 150, 190}
\definecolor{darkgray}{RGB}{45, 52, 54}
\definecolor{lightgray}{RGB}{245, 247, 250}
\definecolor{alertred}{RGB}{214, 48, 49}
\definecolor{succgreen}{RGB}{46, 204, 113}

% Hyperref setup
\hypersetup{
    colorlinks=true,
    linkcolor=cdacblue,
    filecolor=cdacblue,      
    urlcolor=labcyan,
    citecolor=cdacblue,
    pdftitle={TrusMeet: Lab Member and Research Scholar User Manual},
    pdfauthor={Krishna Saha - de.ci.phe.red LAB, IIT Bhilai}
}

\lstdefinelanguage{json}{
    basicstyle=\footnotesize\ttfamily,
    breaklines=true,
    showstringspaces=false,
    frame=lines,
    backgroundcolor=\color{lightgray}
}
\lstset{
    backgroundcolor=\color{lightgray},
    basicstyle=\footnotesize\ttfamily,
    breaklines=true,
    captionpos=b,
    frame=lines,
    rulecolor=\color{cdacblue},
    keywordstyle=\color{cdacblue}\bfseries,
    stringstyle=\color{darkgray},
    commentstyle=\color{gray}\itshape,
    showstringspaces=false
}

% Page style
\pagestyle{fancy}
\fancyhf{}
\fancyhead[L]{\small\textcolor{cdacblue}{\textbf{TrusMeet}} --- Lab Member \& Scholar User Guide}
\fancyhead[R]{\small\textcolor{darkgray}{de.ci.phe.red LAB | IIT Bhilai}}
\fancyfoot[C]{\thepage}
\renewcommand{\headrulewidth}{0.5pt}

% Title formats
\titleformat{\chapter}[display]
  {\normalfont\huge\bfseries\color{cdacblue}}
  {\chaptertitlename\ \thechapter}{10pt}{\Huge}
\titleformat{\section}
  {\normalfont\Large\bfseries\color{cdacblue}}
  {\thesection}{1em}{}
\titleformat{\subsection}
  {\normalfont\large\bfseries\color{darkgray}}
  {\thesubsection}{1em}{}

\tcbuselibrary{skins,breakable}
\newtcolorbox{alertbox}[2][]{
  enhanced,
  breakable,
  colback=white,
  colframe=alertred,
  arc=3mm,
  boxrule=1pt,
  fonttitle=\bfseries\color{white},
  colbacktitle=alertred,
  title=#2,
  #1
}

\newtcolorbox{infobox}[2][]{
  enhanced,
  breakable,
  colback=white,
  colframe=cdacblue,
  arc=3mm,
  boxrule=1pt,
  fonttitle=\bfseries\color{white},
  colbacktitle=cdacblue,
  title=#2,
  #1
}

\newtcolorbox{tipbox}[2][]{
  enhanced,
  breakable,
  colback=white,
  colframe=succgreen,
  arc=3mm,
  boxrule=1pt,
  fonttitle=\bfseries\color{white},
  colbacktitle=succgreen,
  title=#2,
  #1
}

\begin{document}

% Title Page
\begin{titlepage}
    \centering
    \vspace*{1.5cm}
    
    {\Huge\bfseries\color{cdacblue} TrusMeet: Zero-Trust Lab Video Conference Platform\par}
    \vspace{0.8cm}
    {\LARGE\bfseries\color{darkgray} Lab Member \& Research Scholar User Manual\par}
    \vspace{0.4cm}
    {\large Practical Guide to Hardware Token Authentication, Meeting Collaboration, and Secure Participation\par}
    
    \vspace{2.5cm}
    \begin{tcolorbox}[colback=lightgray,colframe=cdacblue,width=0.9\textwidth,arc=2mm]
    \centering
    \textbf{Target Audience:} PhD Scholars, Postdocs, Research Assistants, and Project Students\\
    \textbf{Laboratory:} de.ci.phe.red LAB (Cryptography \& Security)\\
    \textbf{Hardware Device:} C-DAC TrusToken USB Security Key (IS-16695 / ISO-7816)
    \end{tcolorbox}
    
    \vspace{2.5cm}
    \begin{tabular}{ll}
    \textbf{Prepared By:} & Krishna Saha (AI Security) \\
    \textbf{Laboratory:} & de.ci.phe.red LAB (Cryptography \& Information Security) \\
    \textbf{Affiliation:} & Department of CSE, Indian Institute of Technology Bhilai \\
    \textbf{Principal Investigator:} & Dr. Dhiman Saha \\
    \textbf{Version:} & 1.0 (Lab Edition) \quad | \quad \textbf{Date:} November 2026
    \end{tabular}
    
    \vfill
\end{titlepage}

\tableofcontents
\newpage

% CHAPTER 1
\chapter{Introduction \& System Requirements}

\section{Welcome to TrusMeet}
Welcome to \textbf{TrusMeet}, the official high-assurance research video conferencing platform of the \textbf{de.ci.phe.red LAB} at IIT Bhilai. 

Unlike generic meeting tools, TrusMeet is rooted in cryptographic hardware authentication. Each lab member is issued a personal **C-DAC TrusToken** USB dongle containing their unique digital identity certificate. This guarantees that only authorized lab researchers can join sensitive research reviews, cryptanalysis syncs, and thesis evaluations.

\section{System \& Hardware Requirements}
You can join TrusMeet meetings from your personal laptop, lab workstation, or mobile phone.

\begin{table}[h!]
\centering
\small
\begin{tabularx}{\textwidth}{p{3.5cm}X}
\toprule
\textbf{Component} & \textbf{Requirement / Recommendation} \\
\midrule
\textbf{Workstation OS} & Ubuntu Linux 22.04/24.04 LTS, Debian 12, macOS, or Windows 10/11. \\
\textbf{Mobile Devices} & Android (Chrome/Firefox) or iOS (Safari/Chrome). \\
\textbf{Supported Browsers} & Google Chrome ($\ge 105$), Chromium, Microsoft Edge, Brave, Mozilla Firefox ($\ge 102$). \\
\textbf{Hardware Token} & Official C-DAC TrusToken USB Cryptographic Dongle (Model v3.3). \\
\textbf{Peripherals} & Integrated or USB webcam, microphone / headset, and an available USB-A port (or USB-C adapter). \\
\bottomrule
\end{tabularx}
\caption{System Requirements}
\end{table}

\begin{tipbox}{No Software Installation Required on Client}
TrusMeet runs directly inside your web browser. You do not need to install any heavy proprietary desktop clients. All WebRTC streaming and Web Audio processing execute in-browser.
\end{tipbox}

\newpage
% CHAPTER 2
\chapter{Account Creation, Sign-Up \& Authentication}

\section{The Secret Lab Sign-Up Code}
To enforce zero-trust security and prevent unauthorized external internet visitors from creating accounts on the lab server, TrusMeet requires a \textbf{Secret Lab Sign-Up Code}.

\begin{infobox}{Lab Registration Secret}
The sign-up code is provisioned by the Principal Investigator / Lab Administrator. For authorized scholars of \textbf{de.ci.phe.red LAB}, the current active registration code is:
\begin{center}
\Large\texttt{\textbf{710b45f5787e}}
\end{center}
\textit{Keep this code confidential within the lab; do not share it outside authorized research personnel.}
\end{infobox}

\section{Step-by-Step Sign-Up Process (``Create Profile'')}
When accessing TrusMeet for the first time, follow these steps to register:
\begin{enumerate}
    \item \textbf{Access the Platform}: Open the meeting or lobby URL provided by the PI:
    \begin{itemize}
        \item \textbf{Local / Campus Network}: \texttt{https://10.10.1.226:8443} (or \texttt{https://localhost:8443})
        \item \textbf{Mobile Cellular / Remote Internet}: \texttt{https://<TUNNEL-URL>.trycloudflare.com}
    \end{itemize}
    \item \textbf{Select the Sign-Up Tab}: In the authentication panel on the left, click the \textbf{``Create Profile''} tab.
    \item \textbf{Complete the Registration Fields}:
    \begin{itemize}
        \item \textbf{Your Name}: Enter your full professional name (e.g., \textit{Krishna Saha} or \textit{Soumen Jana}, up to 60 characters).
        \item \textbf{Username}: Choose a unique identifier of 3 to 32 characters (letters, numbers, dot, dash, or underscore; e.g., \texttt{ksaha} or \texttt{soumen\_phd}).
        \item \textbf{Password}: Enter a strong master passphrase of \textbf{at least 10 characters}. This password decrypts your personal cryptographic container.
        \item \textbf{Sign-up Code}: Enter the official lab secret: \texttt{\textbf{710b45f5787e}}.
    \end{itemize}
    \item \textbf{Submit Registration}: Click the \textbf{``Create Profile''} button.
\end{enumerate}

\section{What Happens Behind the Scenes (PKI Generation)}
When you click ``Create Profile'', TrusMeet executes an automated cryptographic onboarding protocol:
\begin{enumerate}[noitemsep]
    \item \textbf{Keypair Generation}: The Python PKI bridge generates a cryptographically secure key pair (ECDSA P-256 or RSA).
    \item \textbf{X.509 Certificate Issuance}: An official digital certificate is minted and signed by the \textbf{Lab Root CA} (\texttt{lab\_ca.crt}), establishing your verified identity within the laboratory trust tree.
    \item \textbf{Encrypted Key Storage}: Your private key is encrypted with your chosen password using AES-256 and stored in the secure local SQLite vault (\texttt{trustoken\_users.db}).
    \item \textbf{Automatic Session Activation}: A 256-bit cryptographically secure session token is generated, and you are immediately logged in.
\end{enumerate}

\section{How to Log In (Subsequent Visits)}
For all future meetings, you do not need to register again:
\begin{enumerate}
    \item Navigate to the TrusMeet lobby or click the meeting invitation link.
    \item Under the default \textbf{``Log In''} tab, enter your \textbf{Username} and \textbf{Password}.
    \item Click \textbf{``Log In''}.
    \item If you arrived via a direct meeting link (e.g., \texttt{?room=LAB-DECIPHERED-SYNC}), the application will automatically connect you to the conference room after login!
\end{enumerate}

\section{Connecting Your Physical C-DAC TrusToken}
If you have been issued a physical C-DAC TrusToken USB dongle:
\begin{enumerate}
    \item Insert the token into an available USB port.
    \item Check the status LED on the token casing:
    \begin{itemize}
        \item \textbf{Solid Red Light (Normal / Standby)}: Confirms that the device is receiving 5V DC power from the host USB port and the crypto-container is in secure standby awaiting authentication.
        \item \textbf{Blinking Red Light (Active)}: Indicates that cryptographic operations or APDU commands are actively being exchanged with the host.
    \end{itemize}
    \item In the TrusMeet lobby under your account profile, click \textbf{``Add USB Token''} to bind the physical device to your active session.
\end{enumerate}

\begin{tipbox}{Remote or Mobile Scholars without a USB Token}
Scholars joining from smartphones, tablets, or remote laptops without USB ports can participate fully using their verified \textbf{Software Scholar Profile}. The host will admit you from the waiting room.
\end{tipbox}

\section{Browser Permissions \& Room Admission}
\begin{itemize}
    \item \textbf{SSL Security Bypass (Local IP)}: When connecting over LAN on port 8443, click \textbf{``Advanced''} $\rightarrow$ \textbf{``Proceed to IP (unsafe)''} to accept the local self-signed certificate.
    \item \textbf{Camera \& Microphone Access}: Click \textbf{``Allow''} when prompted by your browser. You can preview and test your video and audio in the right-hand panel before entering.
\end{itemize}

\newpage
% CHAPTER 3
\chapter{In-Meeting Collaboration \& Features}

\section{The Video Conference Interface}
Once inside the room, you will see the modern TrusMeet dark glassmorphic dashboard:
\begin{itemize}
    \item \textbf{Video Grid}: Dynamically arranges participant video feeds.
    \item \textbf{Active Speaker HUD}: The border of the participant currently speaking glows with a neon cyan highlight, driven by real-time audio energy analysis.
    \item \textbf{Role Badges}: Displays the authenticated role of each attendee (e.g., \textit{"Principal Investigator"}, \textit{"PhD Scholar"}, \textit{"Security Researcher"}).
\end{itemize}

\section{Meeting Controls (Bottom Toolbar)}
\begin{table}[h!]
\centering
\small
\begin{tabularx}{\textwidth}{cp{4cm}X}
\toprule
\textbf{Icon} & \textbf{Control Action} & \textbf{Description} \\
\midrule
\textbf{Mic} & Toggle Microphone & Mutes or unmutes your audio. Use when listening to avoid background noise. \\
\textbf{Camera} & Toggle Video & Starts or stops your webcam feed. \\
\textbf{Screen} & Share Screen & Allows sharing entire display, application windows (e.g., VS Code, Terminal), or browser tabs. \\
\textbf{Chat} & Toggle Lab Chat & Opens the side drawer for real-time encrypted messaging and link sharing. \\
\textbf{Phone} & Leave Meeting & Gracefully leaves the room and terminates WebRTC peer connections. \\
\bottomrule
\end{tabularx}
\caption{Toolbar Controls}
\end{table}

\section{Presenting Research Papers \& Code Walkthroughs}
To share your slides or code:
\begin{enumerate}[noitemsep]
    \item Click the \textbf{"Share Screen"} button on the bottom control bar.
    \item Select either \textbf{"Entire Screen"} or \textbf{"Window"} (recommended for PDF presentations or code editors).
    \item If sharing a video with audio, check \textbf{"Share system audio"}.
    \item Click \textbf{"Share"}. Your stream will appear as the primary feed for all participants.
    \item To stop presenting, click the blue \textbf{"Stop Sharing"} bar at the bottom of your screen.
\end{enumerate}

\newpage
% CHAPTER 4
\chapter{The Hardware Presence Rule (Heartbeat)}

\section{Continuous Hardware Liveness Check}
TrusMeet enforces continuous cryptographic liveness:
\begin{itemize}
    \item Every 12--15 seconds, a micro-challenge is sent to your C-DAC TrusToken.
    \item Your token answers the challenge via PKCS\#11 to prove that the authorized researcher is still physically present at the computer.
\end{itemize}

\begin{alertbox}{Important Security Rule: Do NOT Unplug Your Token}
The C-DAC TrusToken USB dongle \textbf{must remain plugged into your computer} for the entire duration of the meeting.
\end{alertbox}

\section{What Happens if the Dongle is Unplugged?}
If the USB token is accidentally or deliberately pulled out:
\begin{enumerate}
    \item The client immediately detects the heartbeat failure.
    \item Your \textbf{microphone is automatically muted} and your \textbf{video stream is paused}.
    \item A security alert banner appears: \textit{"TrusToken disconnected! Audio and video paused for security."}
    \item Other participants in the room are notified that your physical token was removed.
\end{enumerate}

\subsection{How to Restore Access}
\begin{enumerate}[noitemsep]
    \item Re-insert your C-DAC TrusToken into the USB port.
    \item Ensure the red LED returns to steady solid.
    \item Click the \textbf{"Reconnect Token"} button on the prompt to re-authenticate with your PIN.
    \item Your camera and microphone will resume automatically.
\end{enumerate}

\newpage
% CHAPTER 5
\chapter{Reviewing \& Co-Signing Meeting Minutes}

\section{The In-Meeting Minutes Pad}
During lab syncs, the PI or designated scribe updates the \textbf{Lab Meeting Minutes} tab in real time. This pad documents:
\begin{itemize}[noitemsep]
    \item Research experiment progress and cryptanalysis milestones.
    \item Paper submission deadlines and task assignments.
    \item Thesis approvals and equipment requisitions.
\end{itemize}

\section{Cryptographic Sign-Off Workflow}
At the conclusion of the meeting:
\begin{enumerate}
    \item All scholars review the minutes text in the side panel.
    \item The PI locks the minutes and triggers the \textbf{Cryptographic PDF Signing Protocol}.
    \item The document is converted to an ISO 32000-1 PDF and signed with a detached CMS/PKCS\#7 digital signature using the hardware token.
    \item A signed copy (\texttt{minutes\_signed.pdf}) is generated and archived.
\end{enumerate}

\section{Verifying the Signed Minutes on Your Machine}
Any lab member can verify the mathematical authenticity of the meeting minutes:
\begin{lstlisting}[language=bash]
pdfsig minutes_signed.pdf
\end{lstlisting}
The output will display:
\begin{lstlisting}[language=bash]
Signature #1:
  - Signer Certificate Common Name: Dr. Dhiman Saha
  - Signing Time: Verified
  - Digest Algorithm: SHA-256
  - Status: The signature is valid. Document has not been modified.
\end{lstlisting}

\newpage
% CHAPTER 6
\chapter{Frequently Asked Questions \& Troubleshooting}

\section{Common Questions}

\subsection{Why does my token show a solid red light?}
The solid red light is normal. It indicates that the token is receiving power and is in secure standby waiting for PIN authentication.

\subsection{Why can't I access camera/mic when opening \texttt{http://10.10.x.x:8080}?}
Modern browsers disable camera and microphone access on non-secure (plain HTTP) network IP addresses. Always use the **HTTPS** link (\texttt{https://...:8443}) or the Cloudflare tunnel link.

\subsection{What should I do if my PIN fails 3 times?}
To protect against brute-force attacks, C-DAC TrusTokens lock after repeated wrong PIN entries. If you have forgotten your PIN, contact the Lab Administrator (Krishna Saha) to reset it via the administrator PUK.

\subsection{Can I attend using my smartphone?}
Yes! Open the shared Cloudflare tunnel URL in Chrome or Safari on your phone. You can use your scholar emulation profile to participate fully in audio, video, and chat.

\section{Quick Diagnostic Checklist}
\begin{table}[h!]
\centering
\small
\begin{tabularx}{\textwidth}{lX}
\toprule
\textbf{Issue} & \textbf{Quick Fix} \\
\midrule
Audio echoing & Use a headset or earphones to eliminate room acoustic feedback. \\
Video lagging & Check your Wi-Fi signal or switch to a wired Ethernet connection. \\
Token not recognized & Ensure the USB connector is pushed firmly into the port. \\
Accidentally blocked mic & Click the padlock icon in the browser address bar and toggle Microphone to \textbf{"Allow"}. \\
\bottomrule
\end{tabularx}
\caption{Quick Diagnostic Guide}
\end{table}

\vspace{2cm}
\begin{center}
\textbf{--- End of Lab Member \& Scholar User Manual ---}\\
\small \textit{de.ci.phe.red LAB $\cdot$ Department of CSE $\cdot$ Indian Institute of Technology Bhilai}
\end{center}

\end{document}

